\subsubsection{Gaussian Process Regression}

Gaussian Process Regression (GPR) là một phương pháp hồi quy phi tham số
(non-parametric Bayesian approach), trong đó ta không trực tiếp học một hàm ánh
xạ cố định, mà thay vào đó định nghĩa một phân phối trên không gian các hàm.

\paragraph{Khái niệm Gaussian Process}

Một Gaussian Process được định nghĩa như sau:
\begin{equation}
  f(\boldsymbol{x}) \sim \mathcal{GP}(m(\boldsymbol{x}), k(\boldsymbol{x},
  \boldsymbol{x}')),
\end{equation}

trong đó:
\begin{itemize}
  \item $m(\boldsymbol{x})$ là hàm mean (thường giả sử bằng 0)
  \item $k(\boldsymbol{x}, \boldsymbol{x}')$ là hàm kernel mô tả độ tương quan
    giữa hai điểm dữ liệu
\end{itemize}

Với bài toán hồi quy, ta giả sử dữ liệu quan sát:
\begin{equation}
  t_i = f(\boldsymbol{x}_i) + \epsilon, \quad \epsilon \sim \mathcal{N}(0,
  \sigma_n^2).
\end{equation}

\paragraph{Kernel RBF}

Trong thực nghiệm, ta sử dụng Radial Basis Function (RBF) kernel:
\begin{equation}
  k(\boldsymbol{x}, \boldsymbol{x}') = \sigma_f^2
  \exp\left(-\frac{\|\boldsymbol{x} - \boldsymbol{x}'\|^2}{2\ell^2}\right),
\end{equation}

trong đó:
\begin{itemize}
    \item $\ell$ là tham số length-scale, điều khiển độ mượt của hàm
    \item $\sigma_f^2$ là variance, điều khiển biên độ dao động
\end{itemize}

RBF kernel giả định rằng các điểm gần nhau trong không gian input sẽ có giá trị
đầu ra tương tự nhau.

\paragraph{Posterior dự đoán}

Với tập huấn luyện $\boldsymbol{X} = \{\boldsymbol{x}_i\}_{i=1}^n$ và
$\boldsymbol{t} = \{t_i\}_{i=1}^n$, ta xây dựng ma trận kernel:
\begin{equation}
  \boldsymbol{K} = K(\boldsymbol{X}, \boldsymbol{X}) + \sigma_n^2
  \boldsymbol{I}.
\end{equation}

Với một điểm kiểm tra $\boldsymbol{x}_*$, ta có:
\begin{equation}
  \boldsymbol{k}_* = K(\boldsymbol{X}, \boldsymbol{x}_*), \quad
  k_{**} = K(\boldsymbol{x}_*, \boldsymbol{x}_*).
\end{equation}

Khi đó phân phối hậu nghiệm (posterior predictive distribution) là Gaussian:
\begin{equation}
  p(f_* \mid \boldsymbol{x}_*, \boldsymbol{X}, \boldsymbol{t}) =
  \mathcal{N}(\mu_*, \sigma_*^2),
\end{equation}

với:
\begin{gather}
  \mu_* = \boldsymbol{k}_*^T \boldsymbol{K}^{-1} \boldsymbol{t}, \\
  \sigma_*^2 = k_{**} - \boldsymbol{k}_*^T \boldsymbol{K}^{-1}
  \boldsymbol{k}_*.
\end{gather}

\paragraph{Log-Marginal Likelihood}

Các tham số kernel $\boldsymbol{\theta} = \{\ell, \sigma_f, \sigma_n\}$
được học bằng cách tối ưu log-marginal likelihood:
\begin{equation}
  \log p(\boldsymbol{t} \mid \boldsymbol{X}) =
  -\frac{1}{2} \boldsymbol{t}^T \boldsymbol{K}^{-1} \boldsymbol{t}
  - \frac{1}{2} \log |\boldsymbol{K}|
  - \frac{n}{2} \log(2\pi).
\end{equation}

Hàm này gồm ba thành phần:
\begin{itemize}
    \item \textbf{Data fit term}: $\boldsymbol{t}^T \boldsymbol{K}^{-1}
      \boldsymbol{t}$
    \item \textbf{Complexity penalty}: $\log |\boldsymbol{K}|$
    \item \textbf{Normalization constant}
\end{itemize}

\paragraph{Tối ưu tham số bằng Gradient Ascent}

Ta tối ưu tham số kernel bằng gradient ascent: \begin{equation}
  \boldsymbol{\theta} \leftarrow \boldsymbol{\theta} + \eta
  \nabla_{\boldsymbol{\theta}} \log p(\boldsymbol{t} \mid \boldsymbol{X})
\end{equation}

Quy trình huấn luyện: \begin{enumerate}
    \item Khởi tạo tham số $\ell, \sigma_f, \sigma_n$
    \item Tính ma trận kernel $\boldsymbol{K}$
    \item Tính log-marginal likelihood
    \item Tính gradient theo từng tham số
    \item Cập nhật tham số
\end{enumerate}

Quá trình này giúp mô hình tự động học độ mượt và độ nhiễu phù hợp với dữ liệu.

\paragraph{Dự đoán và bất định}

Sau khi huấn luyện, mô hình không chỉ đưa ra giá trị dự đoán mà còn cung cấp độ
bất định: \begin{equation}
  f_* \sim \mathcal{N}(\mu_*, \sigma_*^2).
\end{equation} Trong đó:
\begin{itemize}
    \item $\mu_*$ là giá trị dự đoán trung bình
    \item $\sigma_*^2$ biểu diễn độ không chắc chắn
\end{itemize}
